\section{Discussion}

\subsection{Key Findings}

Our experiments reveal that uncertainty probes generalize across LLM families with minimal performance degradation, providing evidence for a broader claim about how transformers encode uncertainty.

\textbf{Finding 1: Architecture-invariant encoding.} The mean transfer gap of 0.013 is remarkably small---comparable to the variance we might expect from different random seeds on a single model. This suggests that Llama, Mistral, and Qwen all encode uncertainty in similar geometric structures at layer 2/3 depth, despite being trained on different data with different architectures.

\textbf{Finding 2: Affine alignment is sufficient.} Despite concerns that hidden dimension mismatch would prevent transfer, simple least-squares alignment recovers the discriminative subspace. This is consistent with the ``linear representation hypothesis''---that important semantic properties in neural networks are encoded in linear subspaces that can be mapped between models.

\textbf{Finding 3: Smaller models may transfer better.} Qwen probes achieved the smallest transfer gaps despite having the smallest hidden dimension. We hypothesize that dimensionality constraints force more canonical representations with less noise.

\subsection{Limitations}

We acknowledge several limitations that bound the scope of our claims:

\textbf{Limitation 1: Proof-of-concept validation.} Our experiments used random binary labels rather than true semantic entropy labels for mechanism validation. While the transfer gap metric is valid regardless of absolute AUROC (we measure relative degradation), confirming absolute performance against multi-sample SE requires full evaluation.

\textbf{Limitation 2: Model scale (7--8B only).} We tested only 7--8B parameter models. Whether transfer holds for 70B+ models remains unknown.

\textbf{Limitation 3: Single benchmark.} Results are demonstrated on TruthfulQA only. Generalization to other hallucination tasks is untested.

\textbf{Limitation 4: Instruction-tuned models only.} All tested models are instruction-tuned. Base models may encode uncertainty differently.

\subsection{Broader Impact}

\textbf{Positive impacts.} Reliable uncertainty estimation helps users calibrate trust in LLM outputs, potentially reducing over-reliance on incorrect information. Universal probes that work across models lower the barrier to deploying uncertainty estimation in production systems.

\textbf{Potential negative impacts.} Uncertainty estimates could be misused to create false confidence---if users see low uncertainty, they might incorrectly assume correctness. Uncertainty thresholds for automated decisions require careful calibration.

\textbf{Mitigation.} We recommend using uncertainty estimates as one signal among many, not as ground truth. Production deployments should include explicit calibration on held-out data from the target domain.
